% Generated draft: requires mathematical and source review.
\documentclass[11pt]{article}
\usepackage[margin=25mm]{geometry}
\usepackage{fontspec}
\setmainfont{Latin Modern Roman}
\usepackage{amsmath,amssymb,mathtools,unicode-math}
\setmathfont{Latin Modern Math}
\usepackage{enumitem,needspace,xurl,hyperref}
\setlength{\parindent}{0pt}
\setlength{\parskip}{5pt}
\setlength{\emergencystretch}{3em}
\newcommand{\N}{\mathbb{N}}
\newcommand{\Z}{\mathbb{Z}}
\newcommand{\Q}{\mathbb{Q}}
\newcommand{\R}{\mathbb{R}}
\newcommand{\C}{\mathbb{C}}
\newcommand{\F}{\mathbb{F}}
\newcommand{\A}{\mathbb{A}}
\newcommand{\PP}{\mathbb{P}}
\newcommand{\E}{\mathbb{E}}
\newcommand{\eps}{\varepsilon}
\newcommand{\dd}{\,\mathrm{d}}
\begin{document}
\section*{UnsolvedMath ID: 30002084}
\textbf{Source problem code:} catalogue entry
\subsection{Definitions and mathematical statement}
Write spacetime as $\R_t\times\R_x^2$ and let $S^2=\{v\in\R^3:|v|=1\}$. In the usual extrinsic convention, a Schrödinger map is a map $u:I\times\R^2\to S^2$ satisfying
\[
\partial_tu=u\times\Delta_xu,\qquad u(0,x)=u_0(x),
\]
where $I$ is a time interval containing $0$, $\times$ is the cross product in $\R^3$, and $\Delta_x=\partial_{x_1}^2+\partial_{x_2}^2$. This equation makes explicit the conventional interpretation of the term Schrödinger map; the supplied statement does not display the equation or its sign convention. Use the Dirichlet-energy normalization
\[
E(v)=\frac12\int_{\R^2}\bigl(|\partial_{x_1}v|^2+|\partial_{x_2}v|^2\bigr)\,\dd x.
\]
For a smooth map $v$ extending continuously to the one-point compactification $\R^2\cup\{\infty\}\simeq S^2$, its homotopy class is described by its degree, with
\[
\deg(v)=\frac{1}{4\pi}\int_{\R^2}v\cdot(\partial_{x_1}v\times\partial_{x_2}v)\,\dd x.
\]
Thus, in this compactified setting, zero homotopy class means $\deg(v)=0$. Continuous extension at infinity is an explicit sufficient setting for this interpretation, not a boundary condition stated in the supplied record; the intended treatment of more general finite-energy maps is unspecified.

\textbf{Question (permanent catalogue entry).} For every initial map $u_0:\R^2\to S^2$ in the intended admissible class, with zero homotopy class and
\[
E(u_0)<2E(Q_1)=8\pi,
\]
does the Schrödinger-map initial-value problem have global well-posedness and scattering? Global well-posedness asks for existence for all times $t\in\R$, uniqueness in the intended solution class, and continuous dependence on the initial data in the intended topology. Scattering asks for the asymptotic behavior specified by the intended scattering framework.

The strict threshold excludes $E(u_0)=8\pi$. The record supplies the equality $E(Q_1)=4\pi$ but does not define $Q_1$. It also does not specify the initial-data regularity, topology for continuous dependence, uniqueness class, scattering variables, comparison evolution, convergence norm, or whether scattering is required in both time directions. Consequently, the displayed question preserves the supplied target but is not yet a fully specified analytic theorem statement; those choices require source clarification and human mathematical review.
\subsection{Short English statement}
The question asks whether every admissible degree-zero Schrödinger map from two-dimensional space into $S^2$, with initial energy strictly below $8\pi$, evolves globally with uniqueness, continuous dependence, and scattering. The source leaves the precise data class, boundary behavior, and meaning of scattering unspecified. Its supplied status is partially solved; that status is dated curation rather than an independently verified conclusion.
\subsection{Sources}
\begin{itemize}
\item Local UnsolvedMath registry: \texttt{data/unsolvedmath/problems.json}, with \texttt{catalog/problems.json} as the additional-record registry; supplied permanent ID \texttt{30002084}, source problem code \texttt{catalogue entry}, title \emph{Large-Data Scattering for Energy-Critical Schrödinger Maps}, collection \emph{Oberwolfach Reports Open Problems}. No individual proposer is supplied. The supplied proposed year is 2012.
\item Supplied report citation: \emph{Nonlinear Evolution Problems (2013)}, Oberwolfach Reports, DOI 10.4171/owr/2012/26, pp.\ 1563--1637. \url{https://doi.org/10.4171/owr/2012/26}. The supplied background describes Oberwolfach Reports as records of mathematical workshops and their open-problem discussions. The citation's 2013 date and the record's 2012 proposed year are retained without reconciliation.
\item Bejenaru--Ionescu--Kenig--Tataru, \emph{Global Schrödinger maps}. Supplied literature link: \url{https://arxiv.org/abs/0807.0265}.
\item Bejenaru--Ionescu--Kenig--Tataru, \emph{Equivariant Schrödinger maps in two spatial dimensions}. Supplied literature link: \url{https://arxiv.org/abs/1112.6122}.
\item Tataru, \emph{Schrödinger maps and open problems}. Supplied literature link: \url{https://eudml.org/doc/275417}.
\item Literature assessment dated exactly \texttt{2026-08-21}: the supplied curation labels the problem partially solved and reports small critical-data global well-posedness and scattering, together with some low-energy equivariant cases, while describing the large-data problem as unresolved. This is supplied dated curation derived from the OWR record, not an independently verified literature assessment; its time-sensitive status and cited claims require independent checking.
\item The supplied statement status is \texttt{corrected\_verified}. Its verification note reports that the mathematical content was checked against the cited official Oberwolfach report and that notation, context, and question form were normalized without changing the source-level claim. This is a supplied verification assertion; no source visit or independent verification is claimed here.
\end{itemize}
\end{document}
